\documentclass[aspectratio=169,10pt]{beamer}
% Shared preamble for the Algorithms and Programming Concepts slides (Beamer, Metropolis)
\usetheme[progressbar=frametitle,numbering=fraction,background=light,block=fill,sectionpage=none]{metropolis}
\usepackage[T1]{fontenc}
\usepackage{FiraSans}
\usepackage[scaled=.85]{FiraMono}
\usepackage{amsmath,amssymb}
\usefonttheme[onlymath]{serif}
\usepackage{booktabs,array,tabularx}
\usepackage{tikz,pgfplots}
\usepackage{xcolor}
\usepackage{listings}
\usepackage{adjustbox}
\usepackage{fontawesome5}
\hyphenpenalty=10000\exhyphenpenalty=10000
\newenvironment{sfig}[1][0.80]{\begin{adjustbox}{max width=\linewidth,max totalheight=#1\textheight,center}}{\end{adjustbox}}
\usetikzlibrary{arrows.meta,positioning,shapes.geometric,shapes.misc,fit,calc,backgrounds,decorations.pathreplacing,patterns,matrix,chains}
\pgfplotsset{compat=1.17}

\definecolor{tdInk}{HTML}{1D2330}
\definecolor{tdMuted}{HTML}{5D6678}
\definecolor{tdLine}{HTML}{D9D5CA}
\definecolor{tdPaper}{HTML}{FBFAF7}
\definecolor{tdBlue}{HTML}{2457D6}
\definecolor{tdGreen}{HTML}{17855A}
\definecolor{tdAmber}{HTML}{B7791F}
\definecolor{tdRed}{HTML}{C8362F}
\definecolor{tdPurple}{HTML}{7A4FD1}
\setbeamercolor{normal text}{fg=tdInk,bg=tdPaper}
\setbeamercolor{alerted text}{fg=tdRed}
\setbeamercolor{example text}{fg=tdGreen}
\setbeamercolor{frametitle}{bg=tdInk,fg=white}
\setbeamercolor{progress bar}{fg=tdBlue,bg=tdLine}
\setbeamercolor{title separator}{fg=tdBlue}
\setbeamercolor{block title}{bg=tdBlue!12,fg=tdInk}
\setbeamercolor{block body}{bg=tdBlue!5,fg=tdInk}
\setbeamercolor{block title alerted}{bg=tdRed!12,fg=tdRed!80!black}
\setbeamercolor{block body alerted}{bg=tdRed!5,fg=tdInk}
\setbeamercolor{block title example}{bg=tdGreen!14,fg=tdGreen!70!black}
\setbeamercolor{block body example}{bg=tdGreen!6,fg=tdInk}
\setbeamerfont{frametitle}{size=\large,series=\bfseries}
\makeatletter\setlength{\metropolis@frametitle@padding}{1.15ex}\makeatother
\setbeamertemplate{navigation symbols}{}
\setbeamerfont{block title}{size=\small}

% ---- code listings (Python 3.14) ----
\lstdefinestyle{py}{language=Python,
  basicstyle=\ttfamily\fontsize{8.4}{10.1}\selectfont\color{tdInk},
  keywordstyle=\color{tdBlue}\bfseries, commentstyle=\color{tdMuted}\itshape, stringstyle=\color{tdGreen},
  morekeywords={match,case,type,self,True,False,None},
  emph={print,len,range,sorted,min,max,sum,enumerate,zip,heappush,heappop,cache},emphstyle=\color{tdPurple},
  showstringspaces=false, columns=fullflexible, keepspaces=true, frame=single, rulecolor=\color{tdLine},
  backgroundcolor=\color{white}, xleftmargin=2pt, xrightmargin=2pt, aboveskip=2pt, belowskip=0pt, framesep=3pt}
\lstdefinestyle{out}{basicstyle=\ttfamily\fontsize{8}{9.6}\selectfont\color{tdInk}, columns=fullflexible, keepspaces=true,
  frame=single, rulecolor=\color{tdGreen!60}, backgroundcolor=\color{tdGreen!6}, xleftmargin=2pt, xrightmargin=2pt, aboveskip=2pt, belowskip=0pt, framesep=3pt}
% \pycode[size]{snippet} shows snippets/<snippet>.py ; \pyout[size]{snippet} shows snippets/<snippet>.out
\newcommand{\pycode}[2][\relax]{{#1\lstinputlisting[style=py]{snippets/#2.py}}}
\newcommand{\pyout}[2][\relax]{{#1\lstinputlisting[style=out]{snippets/#2.out}}}

% ---- tikz styles ----
\tikzset{>={Stealth[length=2.2mm]},
  card/.style={rounded corners=4pt,draw=#1!70!black,fill=#1!10,align=center,inner sep=5pt,font=\small}, card/.default=tdBlue,
  pill/.style={rounded corners=7pt,draw=none,fill=#1,text=white,inner sep=3pt,font=\scriptsize\bfseries,align=center}, pill/.default=tdBlue,
  arr/.style={->,thick,draw=tdMuted},
  cell/.style={draw=tdInk!70,minimum width=7.5mm,minimum height=7.5mm,inner sep=0pt,font=\small\ttfamily,fill=white},
  tn/.style={circle,draw=tdInk!80,fill=white,minimum size=7mm,inner sep=0pt,font=\small\bfseries}}
\newcommand{\key}[1]{\textbf{\color{tdBlue}#1}}
\newcommand{\bad}[1]{\textbf{\color{tdRed}#1}}
\newcommand{\good}[1]{\textbf{\color{tdGreen!70!black}#1}}
\newcommand{\apcfooter}[1]{\setbeamertemplate{footline}{\hfill\tiny\color{tdMuted}Algorithms and Programming Concepts \,|\, #1 \,|\, BUV\hspace{1.6em}\insertframenumber/\inserttotalframenumber\hspace{0.8em}\vspace{4pt}}}
\newcommand{\apctitle}[2]{\title{#1}\subtitle{#2}\author{Dr.\ Quang-Vinh Dang\\ {\small Senior Lecturer (Assoc.\ Prof)}}\institute{British University Vietnam}\date{Algorithms and Programming Concepts \,|\, Year 1 \,|\, Python 3.14}}

\apcfooter{Sorting}
\apctitle{Sorting}{From simple $O(n^2)$ methods to divide and conquer}

% a row of six cells: \suffixrow{y}{list}{k} colours the last k cells; \prefixrow colours the first k
\newcommand{\suffixrow}[3]{\foreach \v [count=\i from 0] in {#2}{\ifnum\i<\numexpr6-#3\relax\def\fc{white}\else\def\fc{tdGreen!24}\fi
  \node[cell,fill=\fc] at (\i*0.8,#1) {\v};}}
\newcommand{\prefixrow}[3]{\foreach \v [count=\i from 0] in {#2}{\ifnum\i<#3\relax\def\fc{tdGreen!24}\else\def\fc{white}\fi
  \node[cell,fill=\fc] at (\i*0.8,#1) {\v};}}
\newcommand{\mnode}[3]{\node[rounded corners=2pt,draw=tdInk!60,fill=#3,inner sep=2.5pt,font=\scriptsize\ttfamily] at (#1) {#2};}

\begin{document}
\begin{frame}[plain]\titlepage\end{frame}

\begin{frame}{Why sort?}
\begin{columns}[T]
\begin{column}{0.5\textwidth}
{\small
Sorted data makes everything else easier:
\begin{itemize}\setlength\itemsep{3pt}
\item \textbf{search} in $\log n$ steps (binary search);
\item \textbf{rank} and \textbf{top-k}: leaderboards, best-selling items;
\item find \textbf{duplicates}: equal items sit next to each other;
\item \textbf{merge} two sorted lists in one pass;
\item present results in a \textbf{readable order}.
\end{itemize}
Sorting is one of the most studied problems in computer science, and a perfect place to see how a clever idea beats a slow one.}
\end{column}
\begin{column}{0.46\textwidth}
\begin{block}{After this lecture you can}
{\small
\begin{itemize}\setlength\itemsep{2pt}
\item trace \textbf{bubble}, \textbf{insertion}, \textbf{merge} and \textbf{quick} sort;
\item explain why two are $O(n^2)$ and two are $O(n\log n)$;
\item see the difference in \textbf{real timings};
\item use \texttt{sorted} and \texttt{key=} in practice;
\item say what \textbf{stable} means and why it matters.
\end{itemize}}
\end{block}
\end{column}
\end{columns}
\end{frame}

\begin{frame}{Bubble sort: swap neighbours that are out of order}
\begin{columns}[T]
\begin{column}{0.5\textwidth}
\begin{sfig}[0.62]
\begin{tikzpicture}
  \suffixrow{0}{5,2,9,1,7,3}{0}   \node[font=\scriptsize,anchor=east,text=tdMuted] at (-0.6,0) {start};
  \suffixrow{-1.2}{2,5,1,7,3,9}{1} \node[font=\scriptsize,anchor=east,text=tdMuted] at (-0.6,-1.2) {pass 1};
  \suffixrow{-2.4}{2,1,5,3,7,9}{2} \node[font=\scriptsize,anchor=east,text=tdMuted] at (-0.6,-2.4) {pass 2};
  \suffixrow{-3.6}{1,2,3,5,7,9}{3} \node[font=\scriptsize,anchor=east,text=tdMuted] at (-0.6,-3.6) {pass 3};
\end{tikzpicture}
\end{sfig}
{\small Each pass \textbf{bubbles} the largest remaining item to the end (green = in its final place). After pass 4 nothing moves, so we can stop.}
\end{column}
\begin{column}{0.48\textwidth}
\pycode{bubble}
\vspace{2mm}
\pyout{bubble_trace}
\end{column}
\end{columns}
\end{frame}

\begin{frame}{Insertion sort: how you sort a hand of cards}
\begin{columns}[T]
\begin{column}{0.5\textwidth}
\begin{sfig}[0.62]
\begin{tikzpicture}
  \prefixrow{0}{5,2,9,1,7,3}{1}    \node[font=\scriptsize,anchor=east,text=tdMuted] at (-0.6,0) {start};
  \prefixrow{-1.2}{2,5,9,1,7,3}{2} \node[font=\scriptsize,anchor=east,text=tdMuted] at (-0.6,-1.2) {insert 2};
  \prefixrow{-2.4}{2,5,9,1,7,3}{3} \node[font=\scriptsize,anchor=east,text=tdMuted] at (-0.6,-2.4) {insert 9};
  \prefixrow{-3.6}{1,2,5,9,7,3}{4} \node[font=\scriptsize,anchor=east,text=tdMuted] at (-0.6,-3.6) {insert 1};
  \prefixrow{-4.8}{1,2,5,7,9,3}{5} \node[font=\scriptsize,anchor=east,text=tdMuted] at (-0.6,-4.8) {insert 7};
  \prefixrow{-6.0}{1,2,3,5,7,9}{6} \node[font=\scriptsize,anchor=east,text=tdMuted] at (-0.6,-6.0) {insert 3};
\end{tikzpicture}
\end{sfig}
{\small Green prefix = already sorted. Take the next item and slide it left until it fits.}
\end{column}
\begin{column}{0.48\textwidth}
\pycode{insertion}
\vspace{2mm}
{\small
\begin{itemize}\setlength\itemsep{2pt}
\item \textbf{Nearly sorted input}: very few shifts, so about $n$ steps.
\item \textbf{Reversed input}: every item travels all the way left, about $n^2/2$ steps.
\item Great for \textbf{small} lists, so real libraries use it inside faster algorithms.
\end{itemize}}
\end{column}
\end{columns}
\end{frame}

\begin{frame}{Why are these two slow? Counting comparisons}
\begin{columns}[T]
\begin{column}{0.5\textwidth}
{\small
Bubble sort: the passes make $(n-1)+(n-2)+\dots+1$ comparisons:
\[\frac{n(n-1)}{2}\approx\frac{n^2}{2}.\]
Insertion sort has the same order in the worst case.

\vspace{2mm}
\textbf{Double the list $\Rightarrow$ four times the work.}}
\vspace{2mm}

\renewcommand{\arraystretch}{1.2}\small
\begin{tabular}{@{}rr@{}}
\toprule
$n$ & about $n^2/2$ steps\\\midrule
1\,000 & 500\,000\\
100\,000 & 5\,000\,000\,000\\
1\,000\,000 & 500\,000\,000\,000\\
\bottomrule
\end{tabular}
\end{column}
\begin{column}{0.46\textwidth}
\begin{alertblock}{A million items}
{\small At $10^9$ simple steps per second, $5\times10^{11}$ steps take about \textbf{8 minutes} in a fast language; in Python, hours. A smarter algorithm needs only about $n\log_2 n\approx 2\times10^7$ steps: a fraction of a second.}
\end{alertblock}
\vspace{2mm}
{\small The smarter idea is \key{divide and conquer}: split the problem, solve the halves, combine the answers.}
\end{column}
\end{columns}
\end{frame}

\begin{frame}{Merge sort: split in half, sort each half, merge}
\begin{sfig}[0.6]
\begin{tikzpicture}[x=1mm,y=1mm]
  \node[font=\small\bfseries] at (28,31) {1. split until single items};
  \mnode{28,26}{38 27 43 3 9 82 10 5}{white}
  \mnode{14,19}{38 27 43 3}{white}   \mnode{42,19}{9 82 10 5}{white}
  \mnode{7,12}{38 27}{white} \mnode{21,12}{43 3}{white} \mnode{35,12}{9 82}{white} \mnode{49,12}{10 5}{white}
  \draw[arr] (28,24) -- (14,21); \draw[arr] (28,24) -- (42,21);
  \draw[arr] (14,17) -- (7,14); \draw[arr] (14,17) -- (21,14); \draw[arr] (42,17) -- (35,14); \draw[arr] (42,17) -- (49,14);
  \begin{scope}[xshift=82mm]
  \node[font=\small\bfseries] at (28,31) {2. merge sorted halves};
  \mnode{7,26}{27 38}{tdGreen!15} \mnode{21,26}{3 43}{tdGreen!15} \mnode{35,26}{9 82}{tdGreen!15} \mnode{49,26}{5 10}{tdGreen!15}
  \mnode{14,19}{3 27 38 43}{tdGreen!15} \mnode{42,19}{5 9 10 82}{tdGreen!15}
  \mnode{28,12}{3 5 9 10 27 38 43 82}{tdGreen!30}
  \draw[arr] (7,24) -- (14,21); \draw[arr] (21,24) -- (14,21); \draw[arr] (35,24) -- (42,21); \draw[arr] (49,24) -- (42,21);
  \draw[arr] (14,17) -- (28,14); \draw[arr] (42,17) -- (28,14);
  \end{scope}
\end{tikzpicture}
\end{sfig}
\vspace{1mm}
{\small Eight items need $\log_2 8=3$ levels of splitting. At every level the merging touches each item once, so the total work is about $n\log_2 n$.}
\end{frame}

\begin{frame}{Merge sort in Python}
\begin{columns}[T]
\begin{column}{0.5\textwidth}
\pycode{merge}
\end{column}
\begin{column}{0.46\textwidth}
{\small
\textbf{\texttt{merge}} walks two sorted lists with two pointers \texttt{i} and \texttt{j}, always taking the smaller front item. Each item is copied once: $O(n)$.

\vspace{2mm}
\textbf{\texttt{merge\_sort}} is \emph{recursive}: it sorts the left half, the right half, then calls \texttt{merge}. The base case is a list of length $\le 1$.

\vspace{2mm}
\textbf{Cost:} $O(n\log n)$ in \emph{every} case, but it needs $O(n)$ extra memory for the merged lists.

\vspace{2mm}
The \texttt{<=} (not \texttt{<}) in \texttt{merge} keeps equal items in their original order: that makes it \key{stable}.}
\end{column}
\end{columns}
\end{frame}

\begin{frame}{Quick sort: partition around a pivot}
\begin{columns}[T]
\begin{column}{0.5\textwidth}
\begin{sfig}[0.55]
\begin{tikzpicture}[x=1mm,y=1mm]
  \mnode{30,34}{7 2 9 4 5 1 8 6}{white}
  \node[font=\scriptsize,text=tdMuted] at (60,34) {pivot = 5 (middle)};
  \mnode{8,20}{2 4 1}{tdBlue!15} \mnode{30,20}{5}{tdAmber!35} \mnode{52,20}{7 9 8 6}{tdRed!12}
  \draw[arr] (26,32) -- (10,22); \draw[arr] (30,32) -- (30,22); \draw[arr] (34,32) -- (50,22);
  \node[font=\tiny,text=tdBlue] at (8,14) {less than 5}; \node[font=\tiny,text=tdAmber!70!black] at (30,14) {equal: done};
  \node[font=\tiny,text=tdRed] at (52,14) {greater than 5};
  \node[font=\scriptsize] at (30,6) {sort the two outer parts the same way};
\end{tikzpicture}
\end{sfig}
{\small Everything in \texttt{less} belongs left of the pivot and everything in \texttt{more} right of it, so after one pass the pivot is already in its final place.}
\end{column}
\begin{column}{0.48\textwidth}
\pycode{quick}
\vspace{2mm}
{\small
\textbf{Average} $O(n\log n)$. \textbf{Worst case} $O(n^2)$: a bad pivot (e.g.\ always the smallest) splits off just one item. Choosing the \emph{middle} (or a random) item makes that rare.}
\end{column}
\end{columns}
\end{frame}

\begin{frame}{Real timings: sorting $n$ random numbers}
\begin{columns}[T]
\begin{column}{0.46\textwidth}
\pyout[\fontsize{7.4}{8.8}\selectfont]{timing}
\vspace{2mm}
{\small
Doubling $n$: bubble and insertion get about \textbf{4$\times$ slower}; merge and quick only a bit more than \textbf{2$\times$}. At $n=4000$ the gap is already about 100$\times$.}
\vspace{1mm}

{\tiny Milliseconds, best of 3, measured by running \texttt{sorting.py}.}
\end{column}
\begin{column}{0.5\textwidth}
\begin{sfig}[0.7]
\begin{tikzpicture}
\begin{axis}[width=9cm,height=6.4cm,xmode=log,ymode=log,xlabel={list size $n$},ylabel={time (ms)},axis lines=left,legend pos=north west,legend style={font=\scriptsize},tick label style={font=\scriptsize},label style={font=\small}]
  \addplot[tdRed,very thick,mark=*,mark size=1.4pt] table[x=n,y=bubble]{snippets/timing_data.out};
  \addplot[tdAmber,very thick,mark=triangle*,mark size=1.6pt] table[x=n,y=insertion]{snippets/timing_data.out};
  \addplot[tdBlue,very thick,mark=square*,mark size=1.4pt] table[x=n,y=merge]{snippets/timing_data.out};
  \addplot[tdPurple,very thick,mark=diamond*,mark size=1.6pt] table[x=n,y=quick]{snippets/timing_data.out};
  \addplot[tdGreen,very thick,mark=o,mark size=1.4pt] table[x=n,y=builtin]{snippets/timing_data.out};
  \legend{bubble,insertion,merge,quick,sorted()}
\end{axis}
\end{tikzpicture}
\end{sfig}
\end{column}
\end{columns}
\end{frame}

\begin{frame}{In real code: \texttt{sorted()}, \texttt{key=} and stability}
\begin{columns}[T]
\begin{column}{0.54\textwidth}
\pycode{builtin}
\vspace{2mm}
\pyout{builtin}
\vspace{1mm}
{\small Mai and Lan both have 8.5. In both outputs the equal items keep their \textbf{original relative order} (Mai before Lan, then sorted by name): that is a \key{stable} sort.}
\end{column}
\begin{column}{0.42\textwidth}
{\small
Python's \texttt{sorted} and \texttt{list.sort} use \textbf{Timsort}: a merge sort that finds runs that are already in order and uses insertion sort on small pieces.

\vspace{2mm}
\pyout{timsort}
\vspace{2mm}
\textbf{Already sorted} data is sorted in about $O(n)$.

\vspace{2mm}
\textbf{Why stability matters:} sort by grade, then by class: the second sort must not scramble the first.

\vspace{2mm}
Use \texttt{key=} for what to compare and \texttt{reverse=True} for descending order. Don't hand-write a sort in production code.}
\end{column}
\end{columns}
\end{frame}

\begin{frame}{Choosing a sorting algorithm}
\renewcommand{\arraystretch}{1.3}
\footnotesize
\begin{tabular}{@{}>{\raggedright\bfseries}p{25mm}>{\raggedright}p{26mm}>{\raggedright}p{30mm}>{\centering}p{13mm}>{\raggedright\arraybackslash}p{44mm}@{}}
\toprule
Algorithm & Best case & Average / worst & Stable? & Notes\\
\midrule
Bubble sort & $O(n)$ with early exit & $O(n^2)$ & yes & easy to teach, never use in practice\\
Insertion sort & $O(n)$ & $O(n^2)$ & yes & excellent for small or nearly sorted lists\\
Merge sort & $O(n\log n)$ & $O(n\log n)$ & yes & predictable; needs $O(n)$ extra memory\\
Quick sort & $O(n\log n)$ & $O(n\log n)$ / $O(n^2)$ & depends & fast in practice; in-place versions are not stable\\
Timsort (Python) & $O(n)$ & $O(n\log n)$ & yes & adapts to existing order\\
\bottomrule
\end{tabular}
\vspace{4mm}

{\small A known limit: any sort that works by comparing items needs $\Omega(n\log n)$ comparisons in the worst case, so merge sort is optimal in that sense.}
\end{frame}

\begin{frame}{Recap and practice}
\begin{columns}[T]
\begin{column}{0.48\textwidth}
\begin{block}{Take-aways}
{\small
\begin{enumerate}\setlength\itemsep{3pt}
\item \key{Bubble} and \key{insertion} sort compare and swap neighbours: $O(n^2)$.
\item \key{Merge} and \key{quick} sort use \textbf{divide and conquer}: $O(n\log n)$.
\item The gap is huge: ${\sim}100\times$ at $n=4000$, and growing with $n$.
\item A \textbf{stable} sort keeps equal items in their original order.
\item In practice call \texttt{sorted(items, key=...)}.
\end{enumerate}}
\end{block}
\end{column}
\begin{column}{0.48\textwidth}
\begin{block}{Practice}
{\small
\begin{enumerate}\setlength\itemsep{3pt}
\item Trace insertion sort by hand on \texttt{[4, 3, 2, 1]}. How many shifts happen?
\item Add an early-exit flag to \texttt{bubble\_sort}. How fast is it on an already sorted list?
\item Make \texttt{quick\_sort} choose a \emph{random} pivot with \texttt{random.choice}. Does the output change?
\item Sort a list of words by length, then alphabetically for equal lengths, using one \texttt{sorted} call.
\item Time \texttt{merge\_sort} for $n = 10^5$. Is it about $n\log n$?
\end{enumerate}}
\end{block}
\end{column}
\end{columns}
\end{frame}
\end{document}
